\chapter{The Gaussian and Transform Calculus}
\label{ch:iii12}

The Gaussian is the model function of transform calculus: it is smooth, rapidly decaying, and mapped to another Gaussian. Differentiation, scaling, and tensor-product arguments then generate a large transform table.

\section*{Learning goals}
After this chapter, the reader should be able to:
\begin{itemize}
\item derive the Fourier transform of a Gaussian and track normalization constants correctly;
\item use differentiation and multiplication identities to compute transforms of Gaussian-weighted polynomials;
\item apply completion-of-the-square arguments in transform integrals;
\item derive heat-kernel and semigroup formulas from Gaussian transform calculus;
\item compare spatial and frequency widths quantitatively for Gaussian families;
\item construct transform identities by differentiating with respect to parameters when justified;
\end{itemize}
\section*{Conceptual roadmap}
\[
\boxed{\text{structure}\;\longrightarrow\;\text{estimate}\;\longrightarrow\;\text{limit or representation}\;\longrightarrow\;\text{application}.}
\]

\section{Gaussian integrals}
The integral of $e^{-\pi x^2}$ over the line is one under the standard normalization.

\section{Self-reciprocity}
With the $2\pi$ Fourier convention, $e^{-\pi|x|^2}$ is its own Fourier transform.


% BEGIN VOL03-EXPANSION III12-example-01
\begin{example}[Gaussian self-transform]\label{ex:iii12-expand-01}
Let \(G(x)=e^{-\pi x^2}\) and \(F=\widehat G\). Differentiating \(F\) and
integrating by parts gives
\[
F'(\xi)=-2\pi\xi F(\xi).
\]
Hence \(F(\xi)=Ce^{-\pi\xi^2}\). The Gaussian integral gives \(F(0)=1\),
so \(C=1\). Therefore the standard Gaussian is its own Fourier transform.
\end{example}
% END VOL03-EXPANSION III12-example-01
\section{Scaling}
Scaled Gaussians transform to reciprocal-width Gaussians with the appropriate determinant factor.


% BEGIN VOL03-EXPANSION III12-example-02
\begin{example}[Scaled Gaussian]\label{ex:iii12-expand-02}
For \(a>0\),
\[
e^{-\pi a x^2}=G(\sqrt a\,x),
\]
so Fourier scaling yields
\[
\widehat{e^{-\pi a x^2}}(\xi)
=a^{-1/2}e^{-\pi\xi^2/a}.
\]
Spatial narrowing and frequency broadening are reciprocal.
\end{example}
% END VOL03-EXPANSION III12-example-02
\section{Differentiated Gaussians}
Polynomial times Gaussian functions arise by differentiating in space or frequency.

\section{Heat kernels}
Gaussian kernels are the fundamental averaging kernels of the heat equation.


% BEGIN VOL03-EXPANSION III12-example-03
\begin{example}[Heat semigroup]\label{ex:iii12-expand-03}
The heat kernel has Fourier multiplier
\[
e^{-4\pi^2t|\xi|^2}.
\]
Multipliers at times \(s\) and \(t\) multiply to the one at time \(s+t\), so
Fourier inversion gives
\[
G_t*G_s=G_{t+s}.
\]
This is the semigroup property of heat flow.
\end{example}
% END VOL03-EXPANSION III12-example-03
\section{Tensor products}
The $n$-dimensional transform of a product of one-dimensional Gaussians factors coordinatewise.

\section{Hermite structure}
Repeated differentiation of Gaussians produces Hermite-polynomial factors and a transform-invariant basis structure.

\section{Transform tables}
Translation, modulation, scaling, differentiation, and convolution generate many explicit transform pairs from a few seeds.

\section{Core structural results}

\begin{theorem}[Gaussian transform]\label{thm:iii12-01}
Under the convention $\widehat f(\xi)=\int f(x)e^{-2\pi ix\xi}dx$, $e^{-\pi x^2}$ transforms to $e^{-\pi\xi^2}$.
\begin{proof}
Differentiate the transform with respect to $\xi$, integrate by parts to obtain an ODE, solve the ODE, and determine the constant from the Gaussian integral at zero.
\end{proof}
\end{theorem}

\begin{theorem}[Scaling rule]\label{thm:iii12-02}
If $f_a(x)=f(ax)$ with $a\ne0$, then $\widehat{f_a}(\xi)=|a|^{-1}\widehat f(\xi/a)$ in one dimension.
\begin{proof}
Substitute $u=ax$ in the transform integral and account for the Jacobian.
\end{proof}
\end{theorem}

\begin{theorem}[Heat-kernel semigroup]\label{thm:iii12-03}
Gaussian heat kernels satisfy $G_t*G_s=G_{t+s}$.
\begin{proof}
Fourier transformation turns convolution into multiplication, while $\widehat G_t$ is an exponential multiplier whose product at times $t,s$ equals the multiplier at $t+s$; invert the transform.
\end{proof}
\end{theorem}

\section{Worked examples}

\begin{example}[Gaussian integrals diagnostic]\label{ex:iii12-worked-01}
Give a rigorous worked analysis of Gaussian integrals in the setting of this chapter. State the decisive definition or hypothesis and derive the central conclusion. The integral of $e^{-\pi x^2}$ over the line is one under the standard normalization. The decisive point is to use the chapter definition at the correct countable, norm, transform, or distributional level rather than rely on pointwise intuition alone.
\end{example}

\begin{example}[Self-reciprocity diagnostic]\label{ex:iii12-worked-02}
Give a rigorous worked analysis of Self-reciprocity in the setting of this chapter. State the decisive definition or hypothesis and derive the central conclusion. With the $2\pi$ Fourier convention, $e^{-\pi|x|^2}$ is its own Fourier transform. The decisive point is to use the chapter definition at the correct countable, norm, transform, or distributional level rather than rely on pointwise intuition alone.
\end{example}

\begin{example}[Scaling diagnostic]\label{ex:iii12-worked-03}
Give a rigorous worked analysis of Scaling in the setting of this chapter. State the decisive definition or hypothesis and derive the central conclusion. Scaled Gaussians transform to reciprocal-width Gaussians with the appropriate determinant factor. The decisive point is to use the chapter definition at the correct countable, norm, transform, or distributional level rather than rely on pointwise intuition alone.
\end{example}

\begin{example}[Differentiated Gaussians diagnostic]\label{ex:iii12-worked-04}
Give a rigorous worked analysis of Differentiated Gaussians in the setting of this chapter. State the decisive definition or hypothesis and derive the central conclusion. Polynomial times Gaussian functions arise by differentiating in space or frequency. The decisive point is to use the chapter definition at the correct countable, norm, transform, or distributional level rather than rely on pointwise intuition alone.
\end{example}

\section{Solved dossiers}
Each dossier is a canonical solved problem. The provenance ledger records which retained corpus rules guided its topic and which dossiers were added freshly to complete the chapter.

\begin{problem}[Gaussian integrals diagnostic]\label{prob:iii12-dossier-01}
Give a rigorous worked analysis of Gaussian integrals in the setting of this chapter. State the decisive definition or hypothesis and derive the central conclusion.
\end{problem}
\begin{solution}
The integral of $e^{-\pi x^2}$ over the line is one under the standard normalization. The decisive point is to use the chapter definition at the correct countable, norm, transform, or distributional level rather than rely on pointwise intuition alone.
\end{solution}

\begin{problem}[Self-reciprocity diagnostic]\label{prob:iii12-dossier-02}
Give a rigorous worked analysis of Self-reciprocity in the setting of this chapter. State the decisive definition or hypothesis and derive the central conclusion.
\end{problem}
\begin{solution}
With the $2\pi$ Fourier convention, $e^{-\pi|x|^2}$ is its own Fourier transform. The decisive point is to use the chapter definition at the correct countable, norm, transform, or distributional level rather than rely on pointwise intuition alone.
\end{solution}

\begin{problem}[Scaling diagnostic]\label{prob:iii12-dossier-03}
Give a rigorous worked analysis of Scaling in the setting of this chapter. State the decisive definition or hypothesis and derive the central conclusion.
\end{problem}
\begin{solution}
Scaled Gaussians transform to reciprocal-width Gaussians with the appropriate determinant factor. The decisive point is to use the chapter definition at the correct countable, norm, transform, or distributional level rather than rely on pointwise intuition alone.
\end{solution}

\begin{problem}[Differentiated Gaussians diagnostic]\label{prob:iii12-dossier-04}
Give a rigorous worked analysis of Differentiated Gaussians in the setting of this chapter. State the decisive definition or hypothesis and derive the central conclusion.
\end{problem}
\begin{solution}
Polynomial times Gaussian functions arise by differentiating in space or frequency. The decisive point is to use the chapter definition at the correct countable, norm, transform, or distributional level rather than rely on pointwise intuition alone.
\end{solution}

\begin{problem}[Heat kernels diagnostic]\label{prob:iii12-dossier-05}
Give a rigorous worked analysis of Heat kernels in the setting of this chapter. State the decisive definition or hypothesis and derive the central conclusion.
\end{problem}
\begin{solution}
Gaussian kernels are the fundamental averaging kernels of the heat equation. The decisive point is to use the chapter definition at the correct countable, norm, transform, or distributional level rather than rely on pointwise intuition alone.
\end{solution}

\begin{problem}[Tensor products diagnostic]\label{prob:iii12-dossier-06}
Give a rigorous worked analysis of Tensor products in the setting of this chapter. State the decisive definition or hypothesis and derive the central conclusion.
\end{problem}
\begin{solution}
The $n$-dimensional transform of a product of one-dimensional Gaussians factors coordinatewise. The decisive point is to use the chapter definition at the correct countable, norm, transform, or distributional level rather than rely on pointwise intuition alone.
\end{solution}

\begin{problem}[Hermite structure diagnostic]\label{prob:iii12-dossier-07}
Give a rigorous worked analysis of Hermite structure in the setting of this chapter. State the decisive definition or hypothesis and derive the central conclusion.
\end{problem}
\begin{solution}
Repeated differentiation of Gaussians produces Hermite-polynomial factors and a transform-invariant basis structure. The decisive point is to use the chapter definition at the correct countable, norm, transform, or distributional level rather than rely on pointwise intuition alone.
\end{solution}

\begin{problem}[Transform tables diagnostic]\label{prob:iii12-dossier-08}
Give a rigorous worked analysis of Transform tables in the setting of this chapter. State the decisive definition or hypothesis and derive the central conclusion.
\end{problem}
\begin{solution}
Translation, modulation, scaling, differentiation, and convolution generate many explicit transform pairs from a few seeds. The decisive point is to use the chapter definition at the correct countable, norm, transform, or distributional level rather than rely on pointwise intuition alone.
\end{solution}

\begin{problem}[Gaussian transform proof dossier]\label{prob:iii12-dossier-09}
Prove the following structural statement and identify the step where its main hypothesis is used: Under the convention $\widehat f(\xi)=\int f(x)e^{-2\pi ix\xi}dx$, $e^{-\pi x^2}$ transforms to $e^{-\pi\xi^2}$.
\end{problem}
\begin{solution}
Differentiate the transform with respect to $\xi$, integrate by parts to obtain an ODE, solve the ODE, and determine the constant from the Gaussian integral at zero. This proof also explains why the stated hypothesis is structural rather than cosmetic.
\end{solution}

\begin{problem}[Scaling rule proof dossier]\label{prob:iii12-dossier-10}
Prove the following structural statement and identify the step where its main hypothesis is used: If $f_a(x)=f(ax)$ with $a\ne0$, then $\widehat{f_a}(\xi)=|a|^{-1}\widehat f(\xi/a)$ in one dimension.
\end{problem}
\begin{solution}
Substitute $u=ax$ in the transform integral and account for the Jacobian. This proof also explains why the stated hypothesis is structural rather than cosmetic.
\end{solution}

\begin{problem}[Heat-kernel semigroup proof dossier]\label{prob:iii12-dossier-11}
Prove the following structural statement and identify the step where its main hypothesis is used: Gaussian heat kernels satisfy $G_t*G_s=G_{t+s}$.
\end{problem}
\begin{solution}
Fourier transformation turns convolution into multiplication, while $\widehat G_t$ is an exponential multiplier whose product at times $t,s$ equals the multiplier at $t+s$; invert the transform. This proof also explains why the stated hypothesis is structural rather than cosmetic.
\end{solution}

\begin{problem}[Chapter synthesis]\label{prob:iii12-dossier-12}
Connect the main definitions and structural theorems of this chapter into one proof strategy. Explain what object is controlled, what estimate or closure principle is used, and what conclusion becomes available.
\end{problem}
\begin{solution}
The Gaussian is the model function of transform calculus: it is smooth, rapidly decaying, and mapped to another Gaussian. Differentiation, scaling, and tensor-product arguments then generate a large transform table. The recurring strategy is to encode the object in the chapter's natural measurable, normed, Fourier, or distributional structure, establish the controlling estimate, and then pass to the desired limit or representation.
\end{solution}

\section{Exercises with complete solutions}
The shorter exercise layer remains separate from the solved dossiers.

\begin{exercise}\label{exr:iii12-01}
State and apply the central principle behind Gaussian integrals. Give one immediate consequence in the setting of this chapter.
\end{exercise}
\begin{hint}
Complete the square in the exponent and compare the transformed Gaussian with the original Gaussian family.
\end{hint}
\begin{solution}
The integral of $e^{-\pi x^2}$ over the line is one under the standard normalization. As an immediate consequence, the same principle can be inserted into later convergence, approximation, transform, or weak-form arguments whenever its hypotheses are verified.
\end{solution}

\begin{exercise}\label{exr:iii12-02}
State and apply the central principle behind Self-reciprocity. Give one immediate consequence in the setting of this chapter.
\end{exercise}
\begin{hint}
Differentiate the Gaussian transform identity with respect to frequency to generate transforms of \(x^k e^{-ax^2}\).
\end{hint}
\begin{solution}
With the $2\pi$ Fourier convention, $e^{-\pi|x|^2}$ is its own Fourier transform. As an immediate consequence, the same principle can be inserted into later convergence, approximation, transform, or weak-form arguments whenever its hypotheses are verified.
\end{solution}

\begin{exercise}\label{exr:iii12-03}
State and apply the central principle behind Scaling. Give one immediate consequence in the setting of this chapter.
\end{exercise}
\begin{hint}
Multiplication by \(x\) in physical space becomes differentiation in frequency; apply the rule instead of reintegrating from scratch.
\end{hint}
\begin{solution}
Scaled Gaussians transform to reciprocal-width Gaussians with the appropriate determinant factor. As an immediate consequence, the same principle can be inserted into later convergence, approximation, transform, or weak-form arguments whenever its hypotheses are verified.
\end{solution}

\begin{exercise}\label{exr:iii12-04}
State and apply the central principle behind Differentiated Gaussians. Give one immediate consequence in the setting of this chapter.
\end{exercise}
\begin{hint}
Differentiate in the Gaussian parameter when justified to produce moments or heat-kernel identities.
\end{hint}
\begin{solution}
Polynomial times Gaussian functions arise by differentiating in space or frequency. As an immediate consequence, the same principle can be inserted into later convergence, approximation, transform, or weak-form arguments whenever its hypotheses are verified.
\end{solution}

\begin{exercise}\label{exr:iii12-05}
State and apply the central principle behind Heat kernels. Give one immediate consequence in the setting of this chapter.
\end{exercise}
\begin{hint}
For the heat kernel, transform the PDE so the time evolution becomes a scalar ODE in frequency.
\end{hint}
\begin{solution}
Gaussian kernels are the fundamental averaging kernels of the heat equation. As an immediate consequence, the same principle can be inserted into later convergence, approximation, transform, or weak-form arguments whenever its hypotheses are verified.
\end{solution}

\begin{exercise}\label{exr:iii12-06}
State and apply the central principle behind Tensor products. Give one immediate consequence in the setting of this chapter.
\end{exercise}
\begin{hint}
Check mass conservation by integrating the Gaussian kernel or evaluating its transform at zero frequency.
\end{hint}
\begin{solution}
The $n$-dimensional transform of a product of one-dimensional Gaussians factors coordinatewise. As an immediate consequence, the same principle can be inserted into later convergence, approximation, transform, or weak-form arguments whenever its hypotheses are verified.
\end{solution}

\begin{exercise}\label{exr:iii12-07}
State and apply the central principle behind Hermite structure. Give one immediate consequence in the setting of this chapter.
\end{exercise}
\begin{hint}
Spatial narrowing corresponds to frequency widening; derive the scaling law quantitatively.
\end{hint}
\begin{solution}
Repeated differentiation of Gaussians produces Hermite-polynomial factors and a transform-invariant basis structure. As an immediate consequence, the same principle can be inserted into later convergence, approximation, transform, or weak-form arguments whenever its hypotheses are verified.
\end{solution}

\begin{exercise}\label{exr:iii12-08}
State and apply the central principle behind Transform tables. Give one immediate consequence in the setting of this chapter.
\end{exercise}
\begin{hint}
Normalize constants by testing the formula at a parameter value whose integral is already known.
\end{hint}
\begin{solution}
Translation, modulation, scaling, differentiation, and convolution generate many explicit transform pairs from a few seeds. As an immediate consequence, the same principle can be inserted into later convergence, approximation, transform, or weak-form arguments whenever its hypotheses are verified.
\end{solution}


% BEGIN VOL03-EXPANSION III12-exercises-01
\section{Graded supplementary exercises}
These exercises extend the preserved foundational set and are graded by mathematical role.

\subsection*{Standard computations}

\begin{exercise}[Gaussian integral]\label{exr:iii12-expand-01}
What is the normalized Gaussian integral?
\end{exercise}
\begin{hint}
Use the chapter normalization.
\end{hint}
\begin{solution}
Integral of exp(-pi x^2) is 1.
\end{solution}

\begin{exercise}[Self transform]\label{exr:iii12-expand-02}
What is the transform of exp(-pi x^2)?
\end{exercise}
\begin{hint}
Use self-reciprocity.
\end{hint}
\begin{solution}
It is exp(-pi xi^2).
\end{solution}

\begin{exercise}[Scaled Gaussian]\label{exr:iii12-expand-03}
Transform exp(-4 pi x^2).
\end{exercise}
\begin{hint}
Use a=4.
\end{hint}
\begin{solution}
One half times exp(-pi xi^2/4).
\end{solution}

\begin{exercise}[Derivative]\label{exr:iii12-expand-04}
Differentiate exp(-pi x^2).
\end{exercise}
\begin{hint}
Use chain rule.
\end{hint}
\begin{solution}
-2 pi x exp(-pi x^2).
\end{solution}

\begin{exercise}[Heat multiplier]\label{exr:iii12-expand-05}
State the heat multiplier.
\end{exercise}
\begin{hint}
Recall transform solution.
\end{hint}
\begin{solution}
exp(-4 pi^2 t |xi|^2).
\end{solution}

\subsection*{Proofs}

\begin{exercise}[Gaussian ODE]\label{exr:iii12-expand-06}
Derive the ODE for the Gaussian transform.
\end{exercise}
\begin{hint}
Differentiate in frequency and integrate by parts.
\end{hint}
\begin{solution}
F'=-2 pi xi F.
\end{solution}

\begin{exercise}[Scaled formula]\label{exr:iii12-expand-07}
Prove the scaled transform formula.
\end{exercise}
\begin{hint}
Apply Fourier scaling to G(sqrt(a)x).
\end{hint}
\begin{solution}
It gives a^(-1/2) G(xi/sqrt(a)).
\end{solution}

\begin{exercise}[Semigroup proof]\label{exr:iii12-expand-08}
Prove the heat-kernel semigroup.
\end{exercise}
\begin{hint}
Transform convolution.
\end{hint}
\begin{solution}
The exponential multipliers multiply and times add.
\end{solution}

\begin{exercise}[Tensor product]\label{exr:iii12-expand-09}
Why does the multidimensional Gaussian transform factor?
\end{exercise}
\begin{hint}
Both Gaussian and kernel split by coordinates.
\end{hint}
\begin{solution}
Fubini turns the integral into a product of one-dimensional transforms.
\end{solution}

\subsection*{Counterexamples and hypothesis testing}

\begin{exercise}[Normalization]\label{exr:iii12-expand-10}
Why do Gaussian constants depend on Fourier convention?
\end{exercise}
\begin{hint}
The placement of 2 pi changes scaling.
\end{hint}
\begin{solution}
Self-reciprocity in this exact form is convention-dependent.
\end{solution}

\begin{exercise}[Differentiate under integral]\label{exr:iii12-expand-11}
Why must the step be justified?
\end{exercise}
\begin{hint}
Need integrable control on moments.
\end{hint}
\begin{solution}
Without domination the formal derivative may be invalid.
\end{solution}

\begin{exercise}[Width]\label{exr:iii12-expand-12}
Why is every scaled Gaussian not literally self-identical under transform?
\end{exercise}
\begin{hint}
Scaling is reciprocal.
\end{hint}
\begin{solution}
Only the normalized unit-width case has the same width.
\end{solution}

\subsection*{Applications and investigations}

\begin{exercise}[Heat smoothing]\label{exr:iii12-expand-13}
Why does heat flow suppress high frequencies?
\end{exercise}
\begin{hint}
Inspect the multiplier.
\end{hint}
\begin{solution}
Large frequencies are exponentially damped.
\end{solution}

\begin{exercise}[Gaussian filtering]\label{exr:iii12-expand-14}
Why is repeated Gaussian smoothing still Gaussian?
\end{exercise}
\begin{hint}
Use the semigroup law.
\end{hint}
\begin{solution}
The time/variance parameters add.
\end{solution}

\subsection*{Challenge problems}

\begin{exercise}[Hermite structure]\label{exr:iii12-expand-15}
Why do repeated Gaussian derivatives remain polynomial times Gaussian?
\end{exercise}
\begin{hint}
Induct with product rule.
\end{hint}
\begin{solution}
Each derivative changes the polynomial and adds a linear Gaussian factor.
\end{solution}

\begin{exercise}[Parameter derivative]\label{exr:iii12-expand-16}
How can differentiating in a generate transforms involving x^2 times a Gaussian?
\end{exercise}
\begin{hint}
Differentiate the Gaussian family before and after transform.
\end{hint}
\begin{solution}
The a derivative inserts a factor proportional to -x^2, generating the transform identity.
\end{solution}
% END VOL03-EXPANSION III12-exercises-01
\section*{Chapter summary}
The chapter now supplies a canonical theorem-and-dossier layer for subsequent measure, Fourier, distributional, Sobolev, and PDE arguments.
